\chapter{外汇市场与跨境金融}
\label{chap:Foreign_Exchange_Markets_and_Cross_Border_Finance}

\begin{openingcase}{外币应收账款价值多少？}

一家中国制造企业向境外客户出售设备，金额为~1000~万美元，三个月后付款。按当前每美元兑人民币~7.20~元的即期汇率，应收账款价值为人民币~7200~万元。但企业三个月后才能收到美元，届时即期汇率可能不同。

企业可以保留未套期保值的外汇风险暴露，在收款时兑换美元；也可以远期卖出未来美元收入，今天就锁定汇率。销售契约决定美元现金流，外汇市场决定其人民币价值，以及汇率风险如何分配。

\end{openingcase}

\begin{tcolorbox}[
  enhanced,
  colback=FinanceBlue!6,
  colframe=FinanceBlue,
  boxrule=0.6pt,
  sharp corners,
  left=8pt,
  right=8pt,
  top=7pt,
  bottom=7pt,
  before skip=1em,
  after skip=1.6em,
  title=\textbf{核心问题},
  coltitle=white,
  colbacktitle=FinanceBlue,
  fonttitle=\bfseries,
  attach boxed title to top left={xshift=6pt,yshift=-2mm},
  boxed title style={
    boxrule=0pt,
    sharp corners,
    left=5pt,
    right=5pt,
    top=2pt,
    bottom=2pt
  }
]
\begin{enumerate}[
  leftmargin=*,
  label=\textcolor{FinanceBlue}{\textbf{\arabic*.}},
  itemsep=0.55em,
  topsep=0.2em
]

\item 当现金流以不同货币计价时，外汇市场如何与国际贸易、投资和跨境融资相互关联？

\item 通过货币供需关系，现货汇率是如何形成的？资本流动、预期以及市场流动性又如何影响汇率的波动和价格发现呢？

\item 现货汇率和远期汇率是如何通过汇率差异以及覆盖性利息平价来相互关联的？企业和投资者如何利用外汇工具来管理汇率风险呢？

\end{enumerate}

\end{tcolorbox}

\section{外汇市场架构}

汇率是指一种货币相对于另一种货币的价格。外汇市场提供了将跨境支付、投资和资金在不同货币之间进行转换的机制。

\subsection{国际贸易、投资与融资需求}

契约现金流的计价币种与企业或投资者衡量其价值的币种不同时，跨境交易便产生货币风险暴露。计价币种规定契约支付以何种货币表示，估值币种则用于衡量请求权或负债的经济价值。两者不同时，契约金额可能保持固定，其以估值币种衡量的价值却随汇率变化。

设~$S_t$~表示每单位外币对应的本币数量。因此，$S_t$~上升表示本币贬值。

假设一家国内出口商在日期$T$时会收到$F$单位的外币。如果这笔收入按照当时有效的汇率进行兑换，那么其以本国货币表示的价值就是

\begin{equation}
V_{D,T}=S_T V_{F,T}.
\label{eq:ch10-export-receivable-value}
\end{equation}

商业契约固定外币金额，但其本币价值随~$S_T$~变化。较高的~$S_T$~提高外币应收款的本币价值。外币应付款产生相反的净风险暴露，因为较高的~$S_T$~增加履行该义务所需的本币成本。

跨境投资将外国资产的收益与汇率波动结合在一起。假设一位国内投资者以$S_0$的价格将本国货币兑换成外币，然后获得一定的外币收益$r^{*}$，之后又用同样的货币进行反向兑换，此时使用的汇率为$S_1$。那么，本国货币的净收益就是……

\begin{equation}
1+r_d=(1+r^{*})\frac{S_1}{S_0}.
\label{eq:ch10-foreign-investment-return}
\end{equation}

本币回报既取决于外国资产回报，也取决于汇率变化。本币升值可能降低国外投资回报的本币价值，本币贬值则可能提高其本币价值。

跨境借款会使负债方面临货币风险。外币贷款可能具有较低的合同利率，但其本国货币的偿还成本也取决于未来的汇率。因此，在比较国内和外资时，需要同时考虑借款条款和货币风险。

\subsection[市场参与者、交易商与电子交易平台]{市场参与者、交易商以及电子交易平台\ }

外汇市场是一个去中心化的网络，而不是单一的交易所。银行、非银行交易商、企业、资产管理者、对冲基金、中央银行以及其他机构通过交易商网络、电子平台以及双边关系进行交易。

参与者进入市场，是为了执行与贸易、投资和融资交易相关的货币转换。交易商通过报价汇率、提供即时执行服务以及管理由此产生的货币头寸来中介这些交易。

交易商通过中介行为创造了多层次的交易过程。企业客户可以与银行进行交易，而银行则可以将这些头寸与另一客户、另一位交易商或通过电子渠道进行对冲。因此，一笔客户交易可能会引发更多的交易活动，因为货币头寸会在市场中重新分配。

\begin{figure}[H]
\centering

\begin{tikzpicture}[
  >=Latex,
  box/.style={
    draw,
    rounded corners,
    align=center,
    minimum width=3.0cm,
    minimum height=0.9cm,
    inner sep=5pt
  },
  two/.style={
    Latex-Latex,
    line width=0.85pt,
    FinanceBlue
  }
]

\node[box] (corp) at (-4.3,1.7)
  {企业和机构\\跨境借款人};

\node[box] (asset) at (-4.3,-1.7)
  {资产管理者和\\机构投资者};

\node[box] (dealer) at (0,0)
  {银行与\\外汇交易商};

\node[box] (inter) at (4.3,1.7)
  {交易与中介活动\\电子平台};

\node[box] (policy) at (4.3,-1.7)
  {中央银行和\\官方机构};

\draw[two] (corp) -- (dealer);
\draw[two] (asset) -- (dealer);
\draw[two] (dealer) -- (inter);
\draw[two] (dealer) -- (policy);

\end{tikzpicture}

\caption{简化的外汇交易架构}
\label{fig:ch10-fx-market-architecture}
\end{figure}

该图是示意图。一些机构投资者和非银行金融机构也直接访问电子交易平台和其他流动性提供者。中央银行主要通过官方交易参与，包括储备管理和外汇干预。

电子交易通过允许从多个来源获取报价和订单，从而减少了市场许多领域的搜索和执行成本。然而，大额交易、流动性较低的货币对以及定制合约可能更依赖于交易商的关系和双边信贷安排。

外汇交易也产生结算义务。双方交付不同币种，往往通过不同司法辖区的支付系统完成。如果一种货币已交付，而另一种尚未收到，已付款一方就面临结算风险。通过将两种货币的结算相联系，同步交收（PvP）安排可以降低这种风险暴露，但流动性、操作和法律风险仍可能存在。

这些机构安排为汇率的报价以及货币头寸的交易提供了环境。

\subsection{即期报价、市场流动性与价格发现}

即期汇率是指根据市场标准在短期内进行结算时，两种货币之间的兑换价格。不同货币对的结算时间可能有所不同，因此具体的结算日期仍然是交易条款的一部分。

根据这种引用规则，$S_t$ 表示从国内参与者视角直接引用的内容。而相互引用则……

\begin{equation}
\widetilde{S}_t=\frac{1}{S_t},
\label{eq:ch10-reciprocal-quote}
\end{equation}
以单位外币计算单位国内货币。如果汇率是每美元约7.20人民币，那么反向报价大约为$1/7.20=0.1389$以人民币计算的美元价值。

报价惯例决定汇率变动如何解释。采用每单位外币对应本币数量的报价时，$S_t$~上升意味着购买一单位外币需要更多本币。因此，从每美元兑人民币~7.20~元变为~7.40~元，表示人民币对美元贬值。

现货交易的报价通常包含买入价和卖出价。对于以本国货币兑换外币的报价，买入价是交易商购买外币的价格，而卖出价则是交易商出售外币的价格。如果报价是每美元7.1980至7.2020人民币，那么报价的价差就是

\begin{equation*}
\text{保留LaTeX命令、公式、交叉引用和数字。}=7.2020-7.1980=0.0040.
\label{eq:ch10-fx-spread}
\end{equation*}

买卖价差是市场流动性的一个维度。较窄的买卖价差通常意味着交易报价规模的交易 immediate 交易成本较低。流动性还取决于市场深度和价格影响。市场深度指的是当前价格附近可获得的交易量，而价格影响则指的是与较大订单相关的执行价格的变化。这些条件可能因货币对、交易时间、市场状况以及交易规模而有所不同。

交叉汇率可以连接不同货币对之间的价格。例如，如果市场报价显示$S^{\mathrm{CNY/USD}}=7.20$人民币兑换1美元，而$S^{\mathrm{USD/EUR}}=1.10$美元兑换1欧元，那么得出的人民币兑换欧元的交叉汇率就是如此。

\begin{equation*}
S^{\mathrm{CNY/EUR}}
=S^{\mathrm{CNY/USD}}\times
S^{\mathrm{USD/EUR}}
=7.92.
\label{eq:ch10-cross-rate}
\end{equation*}

当直接报价和隐含交叉汇率之间的差距足够大，足以覆盖交易成本和执行限制时，就会产生套利机会。在相关货币对之间进行交易通常可以使这些价格保持紧密一致。

\begin{example}{交叉汇率与三角套利}{fx-triangular-arbitrage}

假设市场的报价为：1美元兑换7.20人民币，1欧元兑换1.10美元。这意味着实际汇率为1欧元兑换7.92人民币。如果直接汇率是1欧元兑换8.00人民币，那么在实际市场中，欧元的价格会比通过美元兑换时更高。

忽略交易成本的话，交易者可以通过美元间接获得欧元，并以较高的直接人民币/欧元汇率出售这些欧元。这样的交易通常可以减少不同货币对之间的价格差异。实际上，这种差异必须足够大，以覆盖买卖价差以及其他执行成本。

\end{example}

汇率决定了交易商给出的报价。当客户和投资者的订单与交易商的报价相互交流时，订单流就能反映出买卖压力的变化，进而影响交易商的库存状况，导致报价的调整和头寸的重新分配。因此，观察到的汇率实际上反映了交易活动与可用流动性以及市场制造能力之间的相互作用。

因此，观察到的汇率变动可能反映了多种力量的共同作用。政策利率的变化、通货膨胀预期、经济状况以及风险感知等因素可能会相互影响，所以汇率的变动取决于这些更广泛的信息如何影响不同货币的预期回报。


\section{汇率度量、制度与动态变化}

货币条件的变化、资本流动、预期和风险都会对货币价值产生压力。这些压力如何体现在市场汇率、官方干预或其他政策调整上，取决于汇率制度。

\subsection{名义汇率、实际汇率与有效汇率}

名义汇率表示一种货币相对于另一种货币的价格，而不考虑国内和外国价格水平的差异。它可以将外币价格转换为本国货币，但本身并不能确定国内商品和外国商品的相对价格。

让 $P_t$ 表示国内价格水平，而 $P_t^*$ 表示国外价格水平。术语 $S_tP_t^*$ 用于将国外价格水平转换为本国货币。将这一金额与国内价格水平进行比较，就可以得到实际汇率：

\begin{equation}
q_t
=
\frac{S_t P_t^{*}}{P_t}.
\label{eq:ch10-real-exchange-rate}
\end{equation}

按这一惯例，$q_t$~上升意味着外国商品相对本国商品变贵，表示本币实际贬值；下降则表示实际升值。

对于比例性的变化，

\begin{equation}
\frac{\Delta q}{q}
\approx
\frac{\Delta S}{S}
+
\pi^{*}
-\pi,
\label{eq:ch10-real-exchange-rate-change}
\end{equation}

其中，$\pi$和$\pi^{*}$分别代表国内和国外的通货膨胀率。因此，实际汇率反映了名义汇率的变动以及通货膨胀的差异。如果国内价格上升的速度远远快于国外价格，那么名义贬值并不一定会导致实际贬值。

\begin{example}{名义汇率与实际汇率变动}{nominal-real-exchange-rate}

假设本国货币汇率从每单位外币7.00单位贬值至7.21单位，贬值幅度约为3\%。在同一时期，国内价格上升了5\%，而国外价格则上升了2\%。

使用近似关系，

$$
\frac{\Delta q}{q}
\approx
3\%+2\%-5\%=0.
$$

名义贬值大致被通货膨胀的差异所抵消，因此实际汇率基本保持不变。

\end{example}

双边汇率反映了两种货币之间的关系。当一个国家与多个合作伙伴进行贸易时，相对于单一货币的变动可能无法反映其货币的整体外部价值。因此，有效的汇率是通过使用一组权重来结合双边汇率得出的，这些权重通常与贸易模式有关。

名义有效汇率概括了加权后的名义货币变动情况。实际有效汇率则加入了相对价格或成本调整因素，这一逻辑同样适用于双边实际汇率，而实际有效汇率则适用于一组贸易伙伴的情况。

因此，当一种货币相对于不同的贸易伙伴发生不同变化时，实际汇率可能会与个别的双边汇率有所不同。然而，这些汇率的解释取决于所使用的权重以及价格或成本指数，因此它们仍然只是对外部竞争力的概括性衡量，而非完整的衡量标准。


\subsection{固定、浮动与管理汇率制度}

不同汇率制度在处理外汇市场压力时，对汇率波动和货币当局政策行动的影响方式有所不同。

在固定或紧密挂钩的汇率制度下，当局试图将汇率维持在一定的平价水平或狭窄的范围内。因此，市场压力更多地通过干预和相关政策调整来加以缓解，不过，如果市场压力持续存在，最终可能需要对汇率安排进行改变。

在浮动汇率制度下，汇率的变动会导致更多的调整。没有固定汇率制度通常意味着货币政策有更大的空间来应对国内状况，尽管仍然可能需要进行干预。

管理型汇率制度结合了这两种调整方式。汇率会根据市场状况进行调节，同时当局仍有权力通过干预或其他操作手段来影响汇率的变动。中国就是一个例子：人民币被正式定义为实行管理型浮动汇率制度，其调整依据的是市场供求情况，同时参考了多种货币篮子。

汇率制度与资本流动性和货币政策独立性之间也存在相互作用。开放经济政策的三重困境意味着，严格固定的汇率、较高的资本流动性以及完全独立的货币政策通常无法同时维持。在汇率固定且资本流动相对开放的情况下，国内货币状况往往会与外部条件更加紧密地联系在一起。

\begin{table}[H]
\centering
\caption{简化的汇率制度与调整机制}
\label{tab:ch10-exchange-rate-regimes}
\begin{tabularx}{\textwidth}{@{}p{2.7cm}X X@{}}
\toprule
\textbf{制度/体制}
& \textbf{主要调整通道}
& \textbf{政策影响} \\
\midrule
已修复或调整了相关设置。
& 干预与政策调整
& 较少的货币政策自主权 \\
已管理好。
& 汇率波动与干预措施
& 中间阶段的政策灵活性 \\
保留LaTeX命令、公式、交叉引用和数字。
& 汇率波动
& 更大的货币政策自主权 \\
\bottomrule
\end{tabularx}
\end{table}

\subsection{资本流动、预期与汇率动态}

投资者会比较国内和外国资产的预期回报，同时考虑利率差异和预期的汇率变动。较高的汇率或资产风险可能会提高投资者要求的回报，而有限的市场流动性则会增加进入或退出市场的成本。资本管制和监管限制可能会进一步限制资金在不同货币之间的流动。

当风险调整后的预期回报相互之间发生变化时，投资者可能会调整他们的投资组合以跨货币进行配置。这些调整会产生资本流动，并改变货币需求，从而影响汇率变动。由于预期汇率变化也会影响预期回报，因此资本流动和汇率是共同决定的。

这些投资组合的调整也受到投资者资金能力和承担风险能力的影响。当杠杆限制收紧或亏损触发保证金要求时，投资者可能会减少外币投资或将资金汇回本国。这种调整可能会加剧汇率波动，并导致短期市场波动，而这种波动并不一定与长期经济前景的变化相匹配。

\begin{tcolorbox}[
  enhanced,
  breakable,
  colback=white,
  colframe=FinanceBlue,
  boxrule=0.7pt,
  sharp corners,
  borderline west={2.2pt}{0pt}{FinanceBlue},
  title=\textbf{金融展望：利率和汇率},
  coltitle=white,
  colbacktitle=FinanceBlue,
  fonttitle=\bfseries,
  before skip=1.2em,
  after skip=1.2em
]

本国利率较高，并不必然导致本币升值。其汇率影响取决于货币政策、通胀与风险的预期路径，以及由此引起的不同币种预期回报变化。

\end{tcolorbox}


\section{平价关系与跨币种估值}

平价关系通过套利和预期，将不同币种的价格和回报联系起来。套利促使不同市场中等价现金流的价值趋于一致，基于预期的关系则取决于对未来价格或汇率的判断。这些关系是在给定假设下的估值基准，而非普遍成立的经验恒等式。

\subsection{一价定律与购买力平价}

一价定律适用于在不同市场出售的相同可交易商品。用$p_{i,t}$表示商品$i$的本国货币价格，用$p_{i,t}^{*}$表示其外币价格。在贸易壁垒和交易成本可以忽略不计的情况下，套利行为通常会在现货汇率转换后使两种价格趋于一致：

\begin{equation}
p_{i,t}=S_t p_{i,t}^{*}.
\label{eq:ch10-law-of-one-price}
\end{equation}

如果国内价格超过转换后的国外价格，且这种差距超过了购买和运输商品的成本，那么交易者就可以在国外购买商品并在国内出售。这种套利行为往往会缩小价格差异。因此，对于那些可交易性强、运输成本低且市场分割程度有限的商品来说，一价定律更有可能成立。

购买力平价将这一逻辑从个别商品扩展到总体价格水平。绝对购买力平价要求可比的商品和服务篮在用共同货币表示时具有相同的成本。考虑到之前定义的实际汇率，绝对购买力平价相当于$q_t=1$。因此，这意味着

\begin{equation}
S_t=\frac{P_t}{P_t^{*}}.
\label{eq:ch10-absolute-ppp}
\end{equation}

实际上，各国的价格指数通常会包含各种可交易与不可交易的商品和服务。消费模式、税收、运输成本以及市场条件的差异都可能导致购买力平价指数无法保持恒定。

相对平价汇率主要关注汇率和价格水平的变化，而非其绝对数值。这一理论基于这样一个假设：实际汇率在一段时间内保持大致不变。正如之前推导出的实际汇率关系所表明的那样，名义汇率会进行调整，以抵消国内与国外之间的通货膨胀差异。根据这里所使用的报价规则，较高的国内通货膨胀意味着本国货币贬值。

汇率并不要求绝对购买力不变。
实际汇率可能会保持在一个不同的水平上，保持稳定。

可以通过一个简单的数值例子来说明这种关系。

\begin{example}{相对购买力平价}
{relative-ppp-example}

假设在未来一年里，国内通货膨胀率预计为5\%，而国外通货膨胀率为2\%。当前的汇率是每单位外币对应7.20个本国货币单位。

在相对购买力平价的基础上，预计本国货币将贬值约3\%。未来汇率的相应购买力平价基准是

\[
S_1
\approx
7.20(1+0.05-0.02)
=
7.416
\approx 7.42.
\]

这个基准反映了预期的通胀差异，而不是对未来汇率的精确预测。

\end{example}

观察到的汇率在较长时间内可能会偏离绝对购买力平价所暗示的水平。运输成本、贸易壁垒、非可交易商品以及市场分割限制了商品市场的套利行为，而价格粘性则减缓了国家价格水平的调整。金融市场的消息也可能使名义汇率的变动速度比商品价格调整的速度快。因此，绝对购买力平价最好被视为汇率水平的长期基准，而相对购买力平价则提供了与通货膨胀差异相关的长期汇率变化的基准。但两者都无法提供可靠的短期汇率预测。


\subsection{抛补利率平价与远期汇率}

抛补利率平价（CIP）是一项无套利条件，使两种投资策略的本币回报相等：直接投资本币资产；或将资金兑换为外币、投资外币资产，并用远期契约对冲汇率风险。投资期限和信用特征可比时，在不存在套利机会的条件下，两种策略必须提供相同的本币回报。

假设有一单位本国货币被投资了一段时间。通过国内投资可以获得$1+r_d$的收益。或者，投资者可以选择按照即期汇率$S_0$将本国货币兑换成外币，然后在国外以$r_f$的利率进行投资，之后再以$F_0$的价格提前出售这些外币收益。这样得到的本国货币支付结果为

\begin{equation*}
\frac{F_0}{S_0}(1+r_f).
\label{eq:ch10-cip-foreign-payoff}
\end{equation*}

如果这两种策略具有相同的到期时间，并且信用和流动性特征相当，那么套利往往会使它们的回报率达到一致：

\begin{equation*}
1+r_d
=
\frac{F_0}{S_0}(1+r_f).
\label{eq:ch10-cip}
\end{equation*}

求解远期汇率后，得到的结果如下：

\begin{equation}
F_0
=S_0\frac{1+r_d}{1+r_f}.
\label{eq:ch10-forward-from-cip}
\end{equation}

对于同一时期内的适度利率水平来说，

\begin{equation}
\frac{F_0-S_0}{S_0}
\approx
r_d-r_f.
\label{eq:ch10-forward-premium-approx}
\end{equation}

根据这种报价规则，当国内利率高于国外利率时，相应的外币远期价格也会高于其现价。

\begin{figure}[H]
\centering
\begin{tikzpicture}[
> =Latex,
box/.style={
draw,
rounded corners,
align=center,
minimum width=3.3cm,
minimum height=0.9cm,
inner sep=5pt
},
flow/.style={-Latex, line width=0.85pt, FinanceBlue}
]
\node[box] (start) at (0,0)
{1个单位的\\国内货币};

\node[box] (dom) at (5.0,1.6)
{国内投资\\在$r_d$处};

\node[box] (for) at (5.0,-1.6)
{转换，投资海外，\\并提前进行套期保值};

\node[box] (domend) at (10.0,1.6)
{$1+r_d$};

\node[box] (forend) at (10.0,-1.6)
{$\dfrac{F_0}{S_0}(1+r_f)$};

\draw[flow] (start) -- (dom);
\draw[flow] (start) -- (for);
\draw[flow] (dom) -- (domend);
\draw[flow] (for) -- (forend);

\end{tikzpicture}
\caption{利率平价中的两条抛补投资路径}
\label{fig:ch10-cip-paths}
\end{figure}

在$T$年的时间内，以连续复合的利率计算，可以用相同的公式来表示这一关系。

\begin{equation}
F_{0,T}
=S_0 e^{(r_d-r_f)T},
\label{eq:ch10-cip-continuous}
\end{equation}
其中，$r_d$和$r_f$分别表示相同到期时间的国内和外国利率。

\begin{example}{由抛补利率平价推导远期汇率}{cip-forward-example}

假设现汇汇率是每单位外币对应7.20个本国货币单位。一年内的本国利率为4\%，而相应的国外利率为2\%。在简化的CIP假设下，

$$
F_0
=
7.20\times\frac{1.04}{1.02}
\approx
7.34.
$$

因此，根据这种报价方式，外币的远期溢价会更高。

\end{example}

实践中，交易成本、抵押品条款、融资条件、交易对手风险、监管和资产负债表约束，都可能使抛补后的回报偏离简单平价基准。此类偏离表明，考虑市场摩擦后，两条融资路径在经济上并非完全相同。

\begin{tcolorbox}[
  enhanced,
  breakable,
  colback=white,
  colframe=FinanceBlue,
  boxrule=0.7pt,
  sharp corners,
  borderline west={2.2pt}{0pt}{FinanceBlue},
  title=\textbf{财务展望：远期利率并非预测结果},
  coltitle=white,
  colbacktitle=FinanceBlue,
  fonttitle=\bfseries,
  before skip=1.2em,
  after skip=1.2em
]

投资者购买外币资产，并提前出售该资产的到期收益，从而锁定了本币回报。如果没有这种套期保值措施，实际获得的回报将取决于到期时的即期汇率。因此，远期汇率实际上是一种约定的未来汇率，而预期的未来即期汇率则反映了投资者的预期以及承担汇率风险所需的补偿。这两者并不一定要一致。

\end{tcolorbox}


\subsection{非抛补利率平价与货币风险溢价}

非抛补利率平价（UIP）比较未对冲汇率风险时不同币种的预期回报。因此，外国存款投资者仍承受未来即期汇率风险。不同于~CIP，UIP~不通过远期契约消除汇率风险，其平价条件涉及预期回报，而非通过契约套期保值锁定的回报。

让$r_d$和$r_f$分别表示在同一投资期限内的国内和外国利率。外国存款的预期本国货币总回报为

\begin{equation*}
(1+r_f)\frac{\mathrm{E}_0[S_1]}{S_0}.
\label{eq:ch10-uip-foreign-return}
\end{equation*}

一个简单的UIP基准测试将这一预期回报与国内总回报等同起来：

\begin{equation}
1+r_d
=
(1+r_f)\frac{\mathrm{E}_0[S_1]}{S_0}.
\label{eq:ch10-uip-exact}
\end{equation}

对于适度的利率和汇率波动，

\begin{equation}
\mathrm{E}_0\left[\frac{S_1-S_0}{S_0}\right]
\approx
r_d-r_f.
\label{eq:ch10-uip-approx}
\end{equation}

本国利率高于外国利率时，UIP~基准意味着预期本币贬值。 

当投资者需要因承担货币风险而获得补偿时，未对冲的外币存款的预期国内货币回报可能会与国内存款的回报有所不同。用$\rho$表示这两种预期回报之间的差值。

\begin{equation}
\rho
=
(1+r_f)\frac{\mathrm{E}_0[S_1]}{S_0}
-
(1+r_d).
\label{eq:ch10-currency-risk-premium}
\end{equation}

因为

\[
\frac{\mathrm{E}_0[S_1]}{S_0}
=
1+
\mathrm{E}_0\left[\frac{S_1-S_0}{S_0}\right],
\]
替换后的结果如下：LaTeX命令、公式、交叉引用以及数字。

\begin{align}
\rho
={}&r_f-r_d
+\mathrm{E}_0\left[\frac{S_1-S_0}{S_0}\right]
+r_f\mathrm{E}_0\left[\frac{S_1-S_0}{S_0}\right].
\end{align}

对于较为温和的利率和汇率变化情况，可以将产品项视为二阶项而忽略不计。由此得到的一阶近似结果为：

\begin{equation}
\rho
\approx
r_f-r_d
+
\mathrm{E}_0\left[\frac{S_1-S_0}{S_0}\right].
\label{eq:ch10-currency-risk-premium-approx}
\end{equation}

重新排序后得到如下结果：LaTeX命令、公式、交叉引用以及数字。

\begin{equation}
\mathrm{E}_0\left[\frac{S_1-S_0}{S_0}\right]
\approx
r_d-r_f+\rho.
\label{eq:ch10-uip-with-risk-premium}
\end{equation}

当$\rho=0$时，两个头寸的预期回报相等，标准UIP基准也就得到了恢复。当$\rho$不为零时，预期汇率变化也反映了承担货币风险所带来的补偿。


\section[货币风险暴露与跨境风险分配]{货币风险暴露与跨境风险分配}

当现金流、资产或负债以与其价值衡量货币不同的货币计价时，就会出现货币风险。汇率变动会导致价值在合同头寸之间重新分配，而套期保值和金融决策则决定了这种风险如何在企业、投资者和金融中介机构之间进行分配。

\subsection{计价币种与契约现金流}

货币单位指的是合同规定的支付单位。以美元计价的债券承诺支付美元款项，而以人民币计价的贷款则需要用人民币进行偿还。货币单位是否会导致企业面临汇率风险，这取决于企业的计价货币以及其收入、资产和其他负债所使用的货币结构。

对于拥有外币负债的公司来说，以本国货币计算的偿还成本会随着即期汇率的变化而变化。本国货币的贬值会增加履行该义务所需的本国货币金额。如果公司的收入主要以本国货币计价，那么外币债务就会形成货币错配。外币收入或资产可以部分抵消这种风险，并提供一种天然的套期保值手段。

合同计价并不能完全涵盖所有汇率风险。企业可以通过选择其结算货币或套期保值合同现金流来减少直接交易风险，同时仍可能因收入、投入成本、竞争环境以及海外业务价值的变化而面临风险。因此，企业的整体货币风险取决于汇率变动对其运营和财务现金流的影响。


\subsection{交易风险暴露}

当以合同形式约定的外币现金流需要在未来日期进行结算时，就会出现交易风险。此时，该现金流的本国货币价值会受到汇率的影响而变得不稳定。这种情况通常出现在应收账款、应付账款以及外币债务偿还过程中。

对方程~\eqref{eq:ch10-export-receivable-value}关于$S_T$求导，同时保持外币应收款项不变，可以得到

\begin{equation}
\frac{\partial V_D,T}{\partial S_T}=V_{F,T}.
\label{eq:ch10-transaction-exposure-sensitivity}
\end{equation}

$S_T$的上升会增加外币应收款的本国货币价值，以及外币应付款的本国货币成本。因此，其经济效应取决于合同中的现金流方向。

由于交易风险通常与特定的金额和结算日期相关联，因此当交易的时间和金额足够可预测时，往往可以与远期合约或其他套期保值合同相结合。

\begin{example}{出口应收账款的交易风险暴露}
{transaction-exposure-exporter}

这位中国出口商将在90天后收到约500万美元的款项。如果结算时的汇率是每美元7.00人民币，那么这笔款项将转换为3500万人民币。但如果汇率达到每美元7.40人民币，那么这笔款项将转换为3700万人民币。

每美元增加0.40人民币的因素，使得应收款项以人民币计算的价值增加了约200万人民币：

\[
\Delta V_{D,T}
=
V_{F,T}\Delta S_T
=
5\times 0.40
=
\text{RMB }2.0\text{ 百万}.
\]

示例展示了固定外币应收款项对汇率的敏感性：较高的$S_T$会使其在本国货币中的价值增加，而较低的$S_T$则会降低这一价值。

\end{example}


\subsection{折算风险暴露}

当需要将以外币表示的财务报表或资产负债表数据转换为本公司的报告货币时，就会涉及到翻译问题。这一点对于拥有海外业务的跨国企业来说尤为重要。

假设一家母公司以本国货币进行报告，而一家外国子公司则使用另一种货币进行账目记录。即使子公司与母公司之间没有相应的现金流变化，汇率的变化仍可能会改变子公司资产、负债、收入和支出的本国货币价值。

因此，折算风险暴露与交易风险暴露不同。交易风险暴露来自将实际结算的契约外币现金流，折算风险暴露则来自将会计金额换算为报告币种。折算损失本身并不产生即时现金支付。

翻译效果仍然可以影响所报告的权益、收益、杠杆率以及其他财务指标。这些指标的重要性取决于公司的资产负债表结构、报告框架，以及与报告中的会计价值相关的任何合同或监管规定。

会计处理方式决定了翻译影响在财务报表中的体现方式。不同的会计框架有着不同的处理规定，因此在进行风险分析时，必须区分实际的外币头寸与其会计记录之间的差异。

\subsection{经济风险暴露}

经济风险暴露有时也称经营风险暴露，指未来经营现金流现值对汇率变化的敏感性。它超出已签约外币应收款和应付款的范围。

即使一家公司的销售额和合同以本国货币计价，它仍然可能面临经济风险。汇率的波动会改变国内产品和国外产品的相对价格、进口投入的成本，以及公司的竞争地位。这些变化可能会影响销售量、利润率以及未来的运营现金流。

经济风险暴露的方向和大小取决于企业收入与成本结构。本币贬值可能提高外币收入价值，但进口投入、外币债务或竞争者反应可能抵消部分影响。因此，并非总能从现有契约的计价币种推断经济风险暴露。


\begin{table}[htbp]
\centering
\caption{三类货币风险暴露}
\label{tab:ch10-currency-exposure-types}

\begin{tabularx}{\textwidth}{@{}p{3.1cm}X X@{}}
\toprule
\textbf{曝光度}
& \textbf{保留LaTeX命令、公式、交叉引用和数字。}
& \textbf{说明性效果} \\
\midrule

交易风险
& 约定的外币现金流
& 结算前的本国货币价值变动 \\

翻译的曝光度&将外币交易转换为报告货币单位&报告的会计数值会随着翻译速率的变化而变化。\\

经济影响
& 汇率对价格、成本、需求和竞争的影响
& 未来现金流现值的变动 \\

\bottomrule
\end{tabularx}
\end{table}

经济风险也会随着企业调整采购、生产、定价和融资方式而发生变化。基于当前运营安排的中短期风险可能会与企业在应对汇率变化后的长期风险有所不同。交易风险、换算风险以及经济风险的主要来源和影响因素在表~\ref{tab:ch10-currency-exposure-types}中有所总结。
这三种风险形式可能会相互重叠。一个国外的业务活动可能会产生换算风险和经济风险，而特定的外币合同则可能导致企业内部的交易风险。


\subsection{套期保值逻辑与汇率风险分配}

套期保值改变头寸对汇率变动的敏感性。其目的通常是抵消已识别的风险暴露，而非预测汇率方向。

考虑一个出口商，其以外币计价的应收款项为$V_{F,T}$，到期日为$T$。如果出口商以汇率$F_{0,T}$将应收款项的一部分$h$进行远期交易，那么相对于即期汇率的结算收益或损失为

\begin{equation}
hV_{F,T}\times(F_{0,T}-S_T).
\label{eq:ch10-export-forward-payoff}
\end{equation}

将预期支付与剩余未对冲应收款的转换结合起来，就可以得到对冲头寸的本国货币价值：

\begin{equation}
V_{D,T}^{H}
=
V_{F,T}\left[(1-h)S_T+hF_{0,T}\right].
\label{eq:ch10-partial-forward-hedge}
\end{equation}

当 $h=0$ 时，应收款项完全暴露于未来的即期汇率之下。而当 $h=1$，且远期合约的金额和到期日与应收款项相匹配时……

\begin{equation}
V_{D,T}^{H}=F_{0,T}V_{F,T}.
\label{eq:ch10-full-forward-hedge}
\end{equation}

因此，套期保值可以用预先确定的汇率来替代部分或全部的不确定的未来转换率。这降低了指定现金流的汇率敏感性，而与基础交易相关的其他风险仍然存在。

\begin{example}{完全与部分远期套期保值}{forward-hedge-example}

这家公司将在六个月内获得约400万美元的收益，而六个月的远期汇率是每美元约7.25人民币。进行全面的套期保值后，其风险将得到锁定。

$$
4\times 7.25)=29.0
$$
百万人民币。

如果该公司选择对冲60\%的应收款项，那么以本地货币表示的价值就会变为……

$$
4\times\left[0.4S_T+0.6\times 7.25\right].
$$

在未来的汇率为每美元7.00人民币的情况下，该收据的价值为2800万人民币；而在汇率为每美元7.50人民币的情况下，其价值则达到2940万人民币。进行部分套期保值虽然可以降低对未来汇率的敏感性，但并不能完全消除这种风险。

\end{example}

同样的逻辑也适用于其他套期保值工具。期货提供标准化的合约，期权会改变汇率结果的分布，而货币互换则可以改变不同货币之间的现金流。企业还可以通过调整收入、成本、资产和负债的货币来创建天然的套期保值措施。在每种情况下，相关的问题都是：套期保值支付如何抵消潜在的暴露风险。

套期保值比率未必等于一。如果基础现金流金额或时点不确定，当实际风险暴露不同于套期保值金额时，完全套期保值可能使企业产生新的货币头寸。衍生品套期保值还可能在收到基础现金流之前，产生抵押品或流动性要求。

通过商业或融资合同本身也可以分配汇率风险。以本国货币开具发票的出口商，其交易风险更多地转移给了外国买家；而以未来收入的货币进行借贷则可以将部分风险转移到负债方。因此，合同的计价方式、运营选择以及金融套期保值为参与者分担汇率风险提供了不同的途径。

\begin{tcolorbox}[
enhanced,
breakable,
colback=white,
colframe=FinanceBlue,
boxrule=0.7pt,
sharp corners,
borderline west={2.2pt}{0pt}{FinanceBlue},
title=\textbf{财务展望：应针对套期保值进行干预，而非对预测进行调整。},
coltitle=white,
colbacktitle=FinanceBlue,
fonttitle=\bfseries,
before skip=1.2em,
after skip=1.2em
]

套期保值头寸的金额、币种和时点与基础风险暴露相对应时，其效果较好。其目的在于改变现金流结果的分布，而非形成对未来汇率的判断。

\end{tcolorbox}


%========================================================
% Summary
%========================================================

\TailModule{summary}{本章小结}

外汇市场将不同货币表示的请求权连接在一起。交易、投资和跨境融资在现金流和计价货币不一致时会产生货币转换的需求。交易商和电子平台在去中心化的市场中重新分配这些流动，而交易和结算安排则影响着流动性和执行效率。

汇率的解释始于报价惯例。在本章中，汇率指的是本国货币与外国货币之间的比率，因此汇率上升意味着本国货币贬值。实际汇率包含了相对价格水平，而有效汇率则反映了与多个贸易伙伴之间的汇率变动情况。不同国家的汇率制度在应对市场压力的方式上有所不同，有的国家通过汇率调整来应对，而有的国家则采取政策手段来干预。

非抛补利率平价在特定的假设条件下提供了基准。一价定律和购买力平价将商品价格与汇率联系起来，这种联系是通过商品市场套利实现的。抛补利率平价则通过抛补回报的比较，将即期汇率、远期汇率和利率联系起来；而非抛补利率平价则将利息差与预期汇率变化联系起来，这些预期汇率变化是基于对预期回报和风险溢价的前提假设而得出的。因此，远期汇率不必等于未来预期的即期汇率。

货币风险敞口源于合同、会计事项、经营现金流以及融资活动。交易风险敞口与以外币计价的现金流相关；转换风险敞口涉及将海外业务转换为报告货币；经济风险敞口则与未来经营现金流对汇率变动的敏感性有关。通过金融对冲和天然对冲措施，可以在现金流的金额、时间以及货币方面与潜在风险敞口保持足够接近的情况下，减少或重新分配这些风险敞口。


%========================================================
% Key Terms
%========================================================

\TailModule{keyterms}{关键术语}

外汇市场（Foreign exchange market）；即期汇率（spot exchange rate）；直接标价（direct quote）；买卖价差（bid--ask spread）；交叉汇率（cross rate）；三角套利（triangular arbitrage）；市场流动性（market liquidity）；名义汇率（nominal exchange rate）；实际汇率（real exchange rate）；名义有效汇率（nominal effective exchange rate）；实际有效汇率（real effective exchange rate）；固定汇率制度（fixed exchange-rate regime）；浮动汇率制度（floating exchange-rate regime）；管理汇率制度（managed exchange-rate regime）；外汇干预（foreign exchange intervention）；一价定律（law of one price）；购买力平价（purchasing power parity）；抛补利率平价（covered interest parity）；远期汇率（forward exchange rate）；非抛补利率平价（uncovered interest parity）；货币风险溢价（currency risk premium）；计价币种（currency denomination）；交易风险暴露（transaction exposure）；折算风险暴露（translation exposure）；经济风险暴露（economic exposure）；远期套期保值（forward hedge）；套期保值比率（hedge ratio）。

%========================================================
% Financial AI Lab
%========================================================

\TailModule{aiprompt}{金融~AI~实验室}

\textbf{实验：对外汇报价进行审计、计算平价，并进行套期保值}

请一个AI系统使用即期汇率、国内和国外的利率以及一年期的远期利率来分析一年的跨货币投资。在进行任何计算之前，要求它先说明汇率的报价方式。

将分析过程分为四个阶段来进行审核。

\begin{enumerate}

\item \textbf{请检查引号的排列方式。} 请确定分子和分母中的货币种类，并确保对升值和贬值的描述与引号中的表述一致。

\item \textbf{确认抛补利率平价是否正确。} 将两个利率设定在相同的到期日和复利基础上。计算由抛补利率平价所暗示的远期汇率，并将其与提供的远期汇率进行比较。在将任何差异视为套利机会之前，需要考虑交易成本、融资和抵押品条件以及其他相关的市场摩擦因素。

\item \textbf{将定价与预测分离开来。} 让人工智能将远期汇率解释为未来即期汇率的预测。检查它是否能够区分覆盖性定价关系与对未来汇率和货币风险溢价的预期。

\item \textbf{对套期保值进行审计。} 给人工智能一个外币应收款项，要求它在没有套期保值的条件下、部分远期套期保值的条件下以及完全远期套期保值的条件下计算现金流。验证远期头寸的方向，并在本国货币升值和贬值的情况下测试结果。

\end{enumerate}


%========================================================
% Discussion Questions
%========================================================

\TailModule{questions}{讨论题}

\begin{enumerate}

\item 为什么外国投资在其本国货币中能够带来正收益，而对于国内投资者来说，在进行货币兑换后却可能带来负收益呢？

\item 为什么拥有浮动汇率的国家仍然可以干预外汇市场？这种干预与维持固定汇率有什么区别？

\item 为什么抛补利率平价提供的套利基准比非抛补利率平价更强？

\item 该公司的出口合同全部以本国货币计价。那么，为什么这家公司仍然会受到汇率波动的影响呢？

\end{enumerate}


%========================================================
% Applied Problems
%========================================================

\TailModule{appprobs}{应用题}

\begin{enumerate}

\item \textbf{汇率、增值率和交叉汇率。}

当前汇率是人民币约7.24元兑换1美元，而美元则兑换欧元约为1.08欧元。

\begin{enumerate}

\item 计算美元/人民币的 reciprocal 报价。

\item 如果人民币对美元的汇率升至7.38，那么根据本章所使用的报价规则，人民币是升值了还是贬值了呢？

\item 请使用原始的报价来计算人民币对欧元的隐含汇率。

\item 解释如果直接报价的人民币/欧元汇率与隐含的交叉汇率相差超过交易成本及其他执行障碍时，可能会产生什么样的交易压力。

\end{enumerate}

\item \textbf{名义汇率和实际汇率。}

在本章所使用的报价制度下，本国货币在一年内贬值了6\%。国内通货膨胀率为4\%，而国外通货膨胀率则降至2\%。

\begin{enumerate}

\item 利用方程~(\ref{eq:ch10-real-exchange-rate-change})中的近似公式，估算实际汇率的变动百分比。

\item 如果国内通货膨胀率为9\%，请重新计算实际的汇率变化，其他数值保持不变。

\item 解释为什么名义折旧并不一定意味着实际折旧。

\end{enumerate}

\item \textbf{抛补利率平价。}

当前汇率为1单位外币对应7.20单位本国货币。一年期本国利率为4.5\%，而相应的国外利率为2.0\%。

\begin{enumerate}

\item 计算由抛补利率平价所暗示的一年远期汇率。

\item 将向前期保费转换为现汇汇率的百分比形式。

\item 假设观察到的远期利率为7.30。在得出存在可交易套利机会的结论之前，还需要哪些额外的信息呢？

\item 请解释为什么CIP所暗示的远期汇率并不必然代表未来的即期汇率。

\end{enumerate}

\item \textbf{未覆盖的回报与货币风险溢价。}

国内一年的利率为5\%，而国外的利率则低至3\%。预计本国货币在一年内将贬值4\%。

\begin{enumerate}

\item 利用方程~\ref{eq:ch10-currency-risk-premium}中的一阶近似公式，计算未对冲的外币存款相对于本国存款的预期超额回报。

\item 在这种假设下，零风险溢价UIP基准测试是否成立呢？

\item 与零风险溢价UIP相对应的预期汇率变化是多少呢？

\item 解释随时间变化的货币风险溢价如何影响利息差异的解读。

\end{enumerate}

\item \textbf{货币风险与套期保值。}

这家国内出口商将在六个月内收到约800万美元。六个月的远期汇率为每美元7.18人民币，该公司选择75\%的套期保值比例。

\begin{enumerate}

\item 如果未来的即期汇率为每美元6.95人民币，请计算对冲后收据的本国货币价值。

\item 如果未来的汇率为每美元7.45人民币，请重新计算数值。

\item 使用相同的两种汇率数据，比较有套期保值措施下的收入范围与没有套期保值措施下的收入范围。

\item 计算在完全套期保值下锁定的本国货币价值。

\item 即使在完整的货币套期保值下，仍存在一个风险。

\end{enumerate}

\end{enumerate}


%========================================================
% Data Application
%========================================================

\TailModule{dataapp}{数据应用}

\textbf{人民币的双边汇率、名义有效汇率和实际有效汇率}

从权威的中国来源获取每月的人民币对美元的汇率，以及从具有 documented methodology 的权威来源获取中国的名义和实际有效汇率指数。

请使用至少五年的数据样本。在适当的情况下，将这一系列数据转换为统一的起始值100。同时绘制出人民币对美元的双边汇率、名义有效汇率指数以及实际有效汇率指数。明确双边汇率的报价原则，并说明每个指数上升的方向。

请指出至少两个时期，在这期间双边汇率和有效指数出现了不同的变化。解释其他贸易伙伴货币的变化如何导致这种差异。然后比较名义有效指数和实际有效指数，并将它们之间的差异与相对价格或成本的变化联系起来。

请仔细解释这些LaTeX命令、公式、交叉引用以及数字的含义。这些元素在描述不同货币之间的汇率变动时起着重要作用，但它们并不能完全反映一个国家的国际竞争力。请至少列举出两个可能影响整体经济分析的因素。

请报告样本周期、频率、指数定义、数据来源机构，以及应用于数据的任何数据处理方法。


%========================================================
% Further Reading
%========================================================

\TailModule{furtherread}{延伸阅读}

\begin{itemize}

\item Dornbusch, R"udiger (1976),
``Expectations and Exchange Rate Dynamics（预期与汇率动态）,''
\emph{Journal of Political Economy}, 84(6), 1161--1176.

\item Meese, Richard A., and Kenneth Rogoff (1983),
``Empirical Exchange Rate Models of the Seventies: Do They Fit Out of Sample?（20世纪70年代的实证汇率模型：样本外表现如何？）''
\emph{Journal of International Economics}, 14(1--2), 3--24.

\item Fama, Eugene F. (1984),
``Forward and Spot Exchange Rates（远期与即期汇率）,''
\emph{Journal of Monetary Economics}, 14(3), 319--338.

\item Obstfeld, Maurice, and Kenneth Rogoff (1995),
``The Mirage of Fixed Exchange Rates（固定汇率的幻象）,''
\emph{Journal of Economic Perspectives}, 9(4), 73--96.

\item Du, Wenxin, Alexander Tepper, and Adrien Verdelhan (2018),
``Deviations from Covered Interest Rate Parity（抛补利率平价的偏离）,''
\emph{The Journal of Finance}, 73(3), 915--957.

\end{itemize}


\ChapterTail{summary,keyterms,aiprompt,questions,appprobs,dataapp,furtherread}
